\subsection{\textit{Modeling} dan Validasi}
\label{subsec:dtl-validasi}

\subsubsection{Skema Validasi}
\todo{Skema validasi untuk DTL (\textit{stratified k-fold} atau \textit{train-test split}), nilai $k$/rasio, dan \textit{seed}. Jika sama dengan kerangka umum (Subbab~\ref{sec:kerangka-umum}), cukup rujuk dan jelaskan perbedaannya saja, misalnya \textit{subsampling} karena waktu komputasi.}

\subsubsection{\textit{Hyperparameter Tuning}}
\todo{Ruang pencarian (\texttt{max\_depth}, \texttt{min\_samples\_leaf}, kriteria \textit{split}, strategi \textit{pruning}, \texttt{confidence}, \texttt{ccp\_alpha}), metode pencarian, dan hasilnya. Sertakan tabel atau grafik skor CV terhadap hiperparameter.}

\begin{table}[H]
	\centering
	\caption{Hasil \textit{tuning} hiperparameter DTL}
	\label{tab:dtl-tuning}
	\begin{tabular}{l l c}
		\toprule
		\textbf{Hiperparameter} & \textbf{Ruang pencarian} & \textbf{Terbaik} \\
		\midrule
		\todo{} & & \\
		\bottomrule
	\end{tabular}
\end{table}

\subsubsection{Model Final}
\todo{Konfigurasi final DTL, skor CV (rata-rata $\pm$ simpangan baku), dan lokasi model tersimpan di \texttt{output/}.}
